% GENERATED FILE. Edit canonical lessons, then run npm run generate:books.
% canonical-source-hash: fnv1a64:e0c56aef49414b51
% canonical-lessons: RU-C224-beryoza, RU-C224-yel, RU-C224-sosna, RU-C224-kust, RU-C224-ovtsa

\chapter{Birch, Spruce, Pine, Bush, Sheep}
\label{ch:ru-birch-spruce-pine-bush-sheep}
\hlchaptermodality{224}

\begin{chapteropening}
\textbf{By the end of this chapter:} \emph{I can say a birch, a spruce, a pine, a bush and a sheep.}

Complete the last lesson of chapter 224: I can say a birch, a spruce, a pine, a bush and a sheep.
\end{chapteropening}

\section[\texorpdfstring{\ru{берёза} (beryóza)}{берёза (beryóza)}]{\ru{берёза} (beryóza) — a birch}

\label{lesson:RU-C224-beryoza}



\begin{quote}
Before the new one: say the Russian for a banana, then the Russian for grapes.
\end{quote}

\subsection*{You'll want to know: \ru{берёза}}
\textbf{\ru{берёза}} (\emph{beryóza}) — \textquotedblleft{}a birch\textquotedblright{}.

\begin{cousinweb}
\textbf{\ru{берёза}} $\leftarrow$ Proto-Slavic \emph{*berza} — inherited. Vasmer lists Lithuanian \emph{béržas}, Sanskrit \emph{bhūrja-}, Old Norse \emph{bjǫrk} and Old High German \emph{birihha}: English \emph{birch} is a cousin, from a root \textquotedblleft{}bright, white\textquotedblright{}.
\end{cousinweb}

\begin{grammarlens}[title={What you've built}]
One more word for this chapter.
\end{grammarlens}

\subsection*{Your turn}
\begin{itemize}
\raggedright
  \item \emph{Say it:} \emph{beryóza}
  \item \emph{Say it:} \emph{beryóza}, once more
  \item \emph{Say it:} say \emph{beryóza}
  \item \emph{Recall:} say the Russian for a banana, then the Russian for grapes, then say \emph{beryóza} again
\end{itemize}

\begin{tcolorbox}[breakable,colback=teal!4,colframe=teal!35!black,title={Before you move on}]
What does \ru{берёза} mean? (\textquotedblleft{}A birch\textquotedblright{}.) Say it once more.
\end{tcolorbox}
\section[\texorpdfstring{\ru{ель} (yel)}{ель (yel)}]{\ru{ель} (yel) — a spruce, a fir tree}

\label{lesson:RU-C224-yel}



\begin{quote}
Before the new one: say the Russian for grapes, then the Russian for a birch.
\end{quote}

\subsection*{You'll want to know: \ru{ель}}
\textbf{\ru{ель}} (\emph{yel}) — \textquotedblleft{}a spruce, a fir tree\textquotedblright{}.

\begin{cousinweb}
\textbf{\ru{ель}} $\leftarrow$ Proto-Slavic \emph{*jedlĭ} — inherited, a cousin of Lithuanian \emph{ẽglė} and Old Prussian \emph{addle}, \textquotedblleft{}fir, spruce\textquotedblright{} (Vasmer).
\end{cousinweb}

\begin{grammarlens}[title={What you've built}]
One more word for this chapter.
\end{grammarlens}

\subsection*{Your turn}
\begin{itemize}
\raggedright
  \item \emph{Say it:} \emph{yel}
  \item \emph{Say it:} \emph{yel}, once more
  \item \emph{Say it:} say \emph{yel}
  \item \emph{Recall:} say the Russian for grapes, then the Russian for a birch, then say \emph{yel} again
\end{itemize}

\begin{tcolorbox}[breakable,colback=teal!4,colframe=teal!35!black,title={Before you move on}]
What does \ru{ель} mean? (\textquotedblleft{}A spruce, a fir tree\textquotedblright{}.) Say it once more.
\end{tcolorbox}
\section[\texorpdfstring{\ru{сосна} (sosná)}{сосна (sosná)}]{\ru{сосна} (sosná) — a pine}

\label{lesson:RU-C224-sosna}



\begin{quote}
Before the new one: say the Russian for a birch, then the Russian for a spruce.
\end{quote}

\subsection*{You'll want to know: \ru{сосна}}
\textbf{\ru{сосна}} (\emph{sosná}) — \textquotedblleft{}a pine\textquotedblright{}.

\begin{cousinweb}
\textbf{\ru{сосна}} $\leftarrow$ Proto-Slavic \emph{*sosna}. Vasmer derives it most likely from a PIE \emph{*\'{k}asnos}, \textquotedblleft{}grey\textquotedblright{}, as in Latin \emph{cānus}, \textquotedblleft{}grey\textquotedblright{}, and Old High German \emph{haso}, \textquotedblleft{}hare\textquotedblright{}; Trubachev, in a note, doubts this.
\end{cousinweb}

\begin{grammarlens}[title={What you've built}]
One more word for this chapter.
\end{grammarlens}

\subsection*{Your turn}
\begin{itemize}
\raggedright
  \item \emph{Say it:} \emph{sosná}
  \item \emph{Say it:} \emph{sosná}, once more
  \item \emph{Say it:} say \emph{sosná}
  \item \emph{Recall:} say the Russian for a birch, then the Russian for a spruce, then say \emph{sosná} again
\end{itemize}

\begin{tcolorbox}[breakable,colback=teal!4,colframe=teal!35!black,title={Before you move on}]
What does \ru{сосна} mean? (\textquotedblleft{}A pine\textquotedblright{}.) Say it once more.
\end{tcolorbox}
\section[\texorpdfstring{\ru{куст} (kust)}{куст (kust)}]{\ru{куст} (kust) — a bush}

\label{lesson:RU-C224-kust}



\begin{quote}
Before the new one: say the Russian for a spruce, then the Russian for a pine.
\end{quote}

\subsection*{You'll want to know: \ru{куст}}
\textbf{\ru{куст}} (\emph{kust}) — \textquotedblleft{}a bush\textquotedblright{}.

\begin{cousinweb}
\textbf{\ru{куст}} $\leftarrow$ Proto-Slavic \emph{*kŭstŭ} — inherited. Vasmer relates it to Lithuanian \emph{kúokštas}, \textquotedblleft{}a bush\textquotedblright{}, and \emph{kùkštas}, \textquotedblleft{}a tuft of straw\textquotedblright{}.
\end{cousinweb}

\begin{grammarlens}[title={What you've built}]
One more word for this chapter.
\end{grammarlens}

\subsection*{Your turn}
\begin{itemize}
\raggedright
  \item \emph{Say it:} \emph{kust}
  \item \emph{Say it:} \emph{kust}, once more
  \item \emph{Say it:} say \emph{kust}
  \item \emph{Recall:} say the Russian for a spruce, then the Russian for a pine, then say \emph{kust} again
\end{itemize}

\begin{tcolorbox}[breakable,colback=teal!4,colframe=teal!35!black,title={Before you move on}]
What does \ru{куст} mean? (\textquotedblleft{}A bush\textquotedblright{}.) Say it once more.
\end{tcolorbox}
\section[\texorpdfstring{\ru{овца} (ovtsá)}{овца (ovtsá)}]{\ru{овца} (ovtsá) — a sheep}

\label{lesson:RU-C224-ovtsa}



\begin{quote}
Before the new one: say the Russian for a pine, then the Russian for a bush.
\end{quote}

\subsection*{You'll want to know: \ru{овца}}
\textbf{\ru{овца}} (\emph{ovtsá}) — \textquotedblleft{}a sheep\textquotedblright{}.

\begin{cousinweb}
\textbf{\ru{овца}} $\leftarrow$ Proto-Slavic \emph{*ovĭca}, a diminutive of the old word for \textquotedblleft{}sheep\textquotedblright{} — inherited. Vasmer lists Sanskrit \emph{ávi-}, Lithuanian \emph{avìs}, Greek \emph{óis} and Latin \emph{ovis}: English \emph{ewe} is a cousin.
\end{cousinweb}

\begin{grammarlens}[title={What you've built}]
One more word for this chapter.
\end{grammarlens}

\subsection*{Your turn}
\begin{itemize}
\raggedright
  \item \emph{Say it:} \emph{ovtsá}
  \item \emph{Say it:} \emph{ovtsá}, once more
  \item \emph{Say it:} say \emph{ovtsá}
  \item \emph{Recall:} say the Russian for a pine, then the Russian for a bush, then say \emph{ovtsá} again
\end{itemize}

\begin{tcolorbox}[breakable,colback=teal!4,colframe=teal!35!black,title={Before you move on}]
What does \ru{овца} mean? (\textquotedblleft{}A sheep\textquotedblright{}.) Say it once more.
\end{tcolorbox}
